% Budget: ca. 1,5 Seiten. Verantwortlich: A (Autonomie), B (Kontext), C (Qualitätssicherung)
% In jedem Unterabschnitt die Wirkung auf Codequalität und Produktivität getrennt benennen.
\section{Ergebnisse}
\label{sec:ergebnisse}

% Überblick: Verteilung der Arbeiten nach Studientyp, Jahr und Venue
\todo{Überblick über die eingeschlossenen Arbeiten}
Tabelle~\ref{tab:uebersicht} ordnet die Arbeiten den drei Dimensionen zu.

\begin{table*}[htbp]
\caption{Übersicht der eingeschlossenen Arbeiten}
\label{tab:uebersicht}
\centering
\begin{tabular}{llllp{6cm}}
\hline
\textbf{Arbeit} & \textbf{Studientyp} & \textbf{Dimension} & \textbf{Wirkung auf} & \textbf{Zentraler Befund} \\
\hline
\todo{Quelle} & \todo{Typ} & \todo{A/K/Q} & \todo{Qualität/Produktivität} & \todo{Befund} \\
\hline
\end{tabular}
\end{table*}

\subsection{Autonomiegrad und menschliche Beteiligung}
\label{sec:ergebnisse-autonomie}
% Wann sind vollautonome Agenten sinnvoll, wann Human-in-the-Loop überlegen?
% Muster Planen, Prüfen, Ausführen
\todo{Befunde zur Qualität}

\todo{Befunde zur Produktivität}

\subsection{Kontext und Steuerung}
\label{sec:ergebnisse-kontext}
% Projektkontextdateien, Spezifikationen, Tests als Leitplanken, standardisierte Werkzeuganbindungen
\todo{Befunde zur Qualität}

\todo{Befunde zur Produktivität}

\subsection{Qualitätssicherung und Fehlermuster}
\label{sec:ergebnisse-qs}
% Automatisierte Tests, Code-Reviews, Multi-Agenten-Ansätze; typische Fehlermuster von Agenten
\todo{Befunde zur Qualität}

\todo{Befunde zur Produktivität}

\todo{Typische Fehlermuster}
